\section{Dataset Penelitian}

Dataset yang digunakan dalam penelitian ini diperoleh melalui proses scraping ulasan Google Reviews menggunakan API dari Apify yang dilakukan secara mandiri. Data dikumpulkan dari beberapa destinasi di Kabupaten Bangkalan yang mewakili kategori wisata alam, wisata religi, dan ruang publik, yaitu Bukit Jaddih, Mangrove Sepulu, Makam Syaikhona Kholil, Aermata Ebhu, dan Lapangan Stadion Bangkalan. Dari proses pengambilan data tersebut diperoleh dataset awal sebanyak 5.250 data dengan 4 atribut (5250 × 4) sebelum dilakukan tahap preprocessing. Dataset ini kemudian digunakan sebagai sumber data utama untuk proses analisis sentimen menggunakan model IndoBERT dan pemodelan topik menggunakan BERTopic guna mengidentifikasi sentimen serta aspek-aspek yang paling banyak dibahas oleh pengunjung pada ulasan destinasi wisata yang diteliti.

\begin{figure}[h]
    \centering
    \includegraphics[width=16cm]{images/Dataset.png}
    \caption{Visualisasi Struktur Dataset Penelitian}
    \label{fig:dataset-penelitian}
\end{figure}